\begin{abstractenv}
    The widespread adoption of Large Language Models (LLMs) via Machine Learning as a Service (MLaaS) platforms has raised significant data privacy concerns. Users are increasingly hesitant to share sensitive information with cloud-based models. This thesis explores the integration of Fully Homomorphic Encryption (FHE) with LLMs to enable privacy-preserving natural language processing. By allowing computations to be performed directly on encrypted data, FHE ensures that neither the user's prompt nor the model's generated response is ever exposed in plaintext to the server. 
    
    This study investigates the architectural adaptations required to execute Transformer-based models within an encrypted domain. Specifically, it addresses the challenges of replacing non-linear activation functions (such as Softmax and GELU) with polynomial approximations suitable for FHE schemes like CKKS. The research evaluates the trade-offs between computational overhead, memory consumption, and inference accuracy, proposing an optimized framework for secure LLM deployment.
\end{abstractenv}